\documentclass[a4paper,11pt]{article}


\usepackage{verbatim}

%- Page layout
%
\pagestyle{empty}
\usepackage[top=30mm,left=30mm,right=30mm,bottom=30mm,nohead]{geometry}


\title{Qpack2}
\author{Patrick Eriksson\\(\texttt{patrick.eriksson@chalmers.se})}
\date{\today}


%--------------------------------------------------------------------------------
\begin{document}


\maketitle

\section{Scope}
\label{sec:scope}
%---
Qpack2 is a package of Matlab functions for performing inversion
work. Only cases where the aim is to retrieve vertical profiles of
atmospheric quantities (1D cases) are handled. The package is also
developed for ground-based observations and inversion of individual
satelitte spectra/scans. This means that, for example, inversion of
satellite scanning sequencies where atmospheric fields to be retrieved
have a latitude or longitude extension (2D or 3D) can not be treated
by Qpack2.

The package should be highly general and flexible for covered
measurements. For example, there is basically no limitations regarding
observation geometry. Further, Qpack2 has been developed to be
suitable for operational inversions, as long as not extreme
calculation speed is required. Some important aspects here are that
batch calculations are allowed and atmospheric a priori profiles can
automatically be extracted from climatology data.

The development of Qpack2 is done using the so called optimal
estimation method as the core inversion methodology. Presently no
other inversion methods are handled, but later additions are possible.



\section{Background}
\label{sec:background}
%---
Qpack2 is part of the more general Atmlab package of Matlab functions.
The computational engine of Qpack2 is ARTS-2. The communication with
ARTS-2 is made through the qarts part of Atmlab. OEM inversions are
perfomred by the function \verb|oem.m| (part of Atmlab).

This can be seen as a direct successor of Qpack, despite all code is
written from scratch. However, for both Qpack versions the
computations are controlled by a set of settings fields packed into a
structure denoted as Q. Qpack was built around ARTS-1, this new Qpack
version is mainly a consequence of that ARTS-2 is the maintained
version. This gave also an oppertunity to make a more stringent
implementation and to add some features. This should hopefully make it
easier to add extensions to Qpack2. Data formats are now more clearly
defined and the ambition level around documentation is higher (but
still relatively modest). A main improvement compared to Qpack2 is
that extraction of a priori data from climatology data is now an
integrated part. On the other hand, the feature of classifying errors
into different categories (0-3 in Qpack) and the associated plotting
features are removed as it was hardly used and has not been used for
operational inversions.



\section{Documentation}
\label{sec:doc}
%---

High level documentation is found in this document. Information on
individual functions is obtained by the standard Matlab help
command. For example:
\begin{verbatim}
>> help qp2_oem
\end{verbatim}
Several of the used datat structures are documented following the
pattern expected of the function \verb|qinfo|. For example, to list
all fields (with documentation text) of \verb|Q| (see below):
\begin{verbatim}
>> qinfo(@qpack2)
\end{verbatim}
The documentation for a specific field is obtained as
\begin{verbatim}
>> qinfo(@qpack2,'F_BACKEND')
\end{verbatim}
Wildcards are allowed when naming fields:
\begin{verbatim}
>> qinfo(@qpack2,'ABS_*')
\end{verbatim}


\section{Main data structures}

\subsection{Forward model and retrieval settings}
\label{sec:Q}
%---
Forward model parameters and retrieval variables are joined into a
single structure, denoted as \verb|Q|. This structure has a fixed set
of fields. That is, the structure must always contain exactly this set
of fields. However, some fields are used only in particularly
circumstances and not all fields must be set.

The fields of \verb|Q| are listed and described in
Appendix~\ref{app:q}. The same information is obtained inside Matlab
through \verb|qinfo|, where the defintionn function is
\verb|qpack2|. The examples in Section~\ref{sec:doc} can be used here.

Many of the fields of \verb|Q| are identical to the corresponding
field of the qarts interface, and additional information can be found
as:
\begin{verbatim}
>> qinfo(@qarts)
\end{verbatim}


\subsection{Measurement data}
\label{sec:Y}
%---
The series of measurements (spectra or scans) to be treated is packed
into a data structure denoted as \verb|Y|. The data fields of this
structure are described in \verb|qp2_y|, and available information can be
obtained as
\begin{verbatim}
>> qinfo(@qp2_y)
\end{verbatim}
The same information is appended to this document in
Appendix~\ref{app:y}.



\section{Standard inversions}
\label{sec:stdinversions}
%---




\verb|Q.F_BACKEND|

\begin{verbatim}
>> Q = q_h2o;
>> Y = get_onsala_Y( 'h2o', date1, date2 );
>> qp2_check_f( Q, Y, 1e3 );
>> Y = rmfield( Y, 'f' );
>> P = qp2_oem( Q, Y );
\end{verbatim}



\appendix

\newpage\noindent{\huge \bf Appendices}

\section{The fields of Q}
\label{app:q}
%
\verbatiminput{qpack2.txt}


\newpage
\section{The fields of Y}
\label{app:y}
%
\dots


\end{document}
